\chapter{Beyond the Weak Sector}
\label{ch:beyond}

\begin{quote}
\textit{This brief epilogue looks ahead to applications of EDC beyond
the weak sector covered in this book.}
\end{quote}

% ════════════════════════════════════════════════════════════
\section{What We Have Established}
% ════════════════════════════════════════════════════════════

This book has shown that the EDC framework, developed in Book 1 for
strong interactions and mass ratios, extends naturally to the weak sector:

\begin{itemize}[nosep]
\item \textbf{Neutron decay}: Junction relaxation through geometric barrier
\item \textbf{V$-$A structure}: Chiral localization from boundary conditions
\item \textbf{Weinberg angle}: $\sin^2\theta_W = 1/4$ from $Z_6$ symmetry
\item \textbf{Three generations}: Angular sectors from $Z_3 \subset Z_6$
\item \textbf{CKM hierarchy}: Wavefunction overlap model
\item \textbf{Neutrino suppression}: Edge mode localization
\end{itemize}

All these emerge from the same 5D membrane geometry---no new parameters
were introduced for the weak sector.

% ════════════════════════════════════════════════════════════
\section{Nuclear Applications}
% ════════════════════════════════════════════════════════════

The next frontier is nuclear physics. EDC provides a natural framework
for understanding:

\subsection{Nuclear Binding}

In EDC, nuclei are not ``bags of nucleons'' but rather extended 5D
structures where multiple junctions share membrane area:

\begin{itemize}[nosep]
\item Binding energy from shared membrane tension
\item Magic numbers from geometric packing
\item Nuclear radius scaling from junction extent
\end{itemize}

\subsection{Nuclear Decay Modes}

The barrier tunneling picture extends to:
\begin{itemize}[nosep]
\item Alpha decay: Emission of a 4-junction cluster
\item Beta decay: Single junction relaxation (covered in this book)
\item Fission: Large-scale membrane splitting
\end{itemize}

\subsection{Preview: Helium-4}

The alpha particle ($^4$He) has special stability in EDC:
\begin{itemize}[nosep]
\item Four junctions form a closed tetrahedral configuration
\item Maximum surface-to-volume efficiency
\item Explains why $^4$He is the ``end product'' of stellar fusion
\end{itemize}

% ════════════════════════════════════════════════════════════
\section{Experimental Tests}
% ════════════════════════════════════════════════════════════

EDC makes predictions that differ from---or are more specific than---the
Standard Model:

\subsection{Precision Tests}

\begin{center}
\begin{tabular}{lll}
\toprule
\textbf{Observable} & \textbf{EDC Prediction} & \textbf{Current Status} \\
\midrule
$\sin^2\theta_W$ & $1/4$ at tree level & 8\% off; radiative corrections expected \\
$m_p/m_e$ & $6\pi^5$ exactly & 0.002\% agreement \\
$M_Z/M_W$ ratio & $2/\sqrt{3}$ & 1.8\% off; within envelope \\
\bottomrule
\end{tabular}
\end{center}

\textbf{Critical test}: If radiative corrections to $\sin^2\theta_W$ can be
calculated in EDC, the prediction becomes falsifiable.

\subsection{Structural Tests}

\begin{itemize}[nosep]
\item \textbf{Fourth generation}: Would falsify $Z_3$ generation model
\item \textbf{Proton decay}: EDC predicts absolute stability (topological protection)
\item \textbf{Neutron lifetime anomaly}: EDC may explain beam vs bottle discrepancy
\end{itemize}

\subsection{Cosmological Tests}

The 5D membrane picture has cosmological implications:
\begin{itemize}[nosep]
\item Dark matter as bulk excitations
\item Dark energy from membrane tension
\item Baryogenesis from topological processes
\end{itemize}

These are speculative but potentially testable.

% ════════════════════════════════════════════════════════════
\section{Philosophical Implications}
% ════════════════════════════════════════════════════════════

EDC changes the status of ``fundamental constants'':

\begin{statusbox}{Constants as Geometry}
In the Standard Model, constants like $\alpha$, $G_F$, and mass ratios
are \textit{inputs}---they ``just are'' what they are.

In EDC, these same quantities are \textit{outputs}---they follow from
5D membrane geometry. This has profound implications:

\begin{itemize}[nosep]
\item Constants could, in principle, be different
\item They are connected by geometric relations
\item An entity with access to 5D could change our physics
\end{itemize}

This doesn't mean such changes will happen---only that EDC provides
a \textit{mechanism} where the Standard Model provides none.
\end{statusbox}

% ════════════════════════════════════════════════════════════
\section{Open Questions}
% ════════════════════════════════════════════════════════════

Major questions remain:

\begin{enumerate}
\item \textbf{Gravity}: How does general relativity emerge from 5D membrane dynamics?

\item \textbf{Quantum mechanics}: Is the wavefunction a 5D object projected to 3D?

\item \textbf{Cosmology}: What is the cosmological history of the membrane?

\item \textbf{Unification}: Can all forces be unified in the 5D picture?
\end{enumerate}

These are beyond the scope of this book but define the research program ahead.

% ════════════════════════════════════════════════════════════
\section{Conclusion}
% ════════════════════════════════════════════════════════════

The EDC framework has now been applied to:
\begin{itemize}[nosep]
\item Mass ratios and fundamental constants (Book 1)
\item Weak interactions and decay processes (Book 2)
\end{itemize}

In both cases, the same geometric principles---membrane tension,
junction topology, $Z_6$ symmetry, frozen projection---produce
quantitative predictions that match observation.

\textbf{No new parameters were introduced for the weak sector.}

This is either an extraordinary coincidence or evidence that the
5D membrane picture captures something real about nature.

The path forward is clear:
\begin{enumerate}[nosep]
\item Close the open problems (Appendix~\ref{app:opr})
\item Apply to nuclear physics
\item Seek experimental discrimination from Standard Model
\item Explore cosmological implications
\end{enumerate}

The weak sector is not the end---it is a waypoint on a larger journey.

\bigskip

\begin{center}
\textit{The 5D causes; the 3D observes.}
\end{center}
